% !TEX root = ../main.tex
\section{Análisis del costo computacional de cada función}

El análisis siguiente reporta dos cosas para cada función: la clase
asintótica (notación Big O) y el costo temporal real que esa clase
esconde. Esta distinción importa porque una función puede clasificarse
como $O(n)$ y aun así estar ejecutando, por ejemplo, $n + \log n$ o $3n$
operaciones reales. La notación Big O descarta constantes y términos de
menor orden, pero esas operaciones se siguen ejecutando en la práctica.
Cuando una función combina dos algoritmos, por ejemplo una búsqueda en el
mapa de posiciones seguida de una lectura en disco, el costo real es la
suma de ambos, no solo el término dominante.

Además del costo temporal, se documenta el costo espacial: cuánta memoria
adicional usa cada función más allá de los datos que ya existen en el
sistema, por ejemplo un acumulador, un arreglo temporal, o ninguna
memoria extra.

\subsection{Módulo de mapas de posiciones}

{\small
\begin{longtable}{p{3.3cm} p{4.3cm} p{2.6cm} p{1.8cm} p{2.9cm}}
\toprule
Función & Algoritmo & Costo real & Big O & Costo espacial \\
\midrule
\endhead

buscarPosicionPorCarne &
    Binaria en zona ordenada, lineal en lista de recién ingresados si
    no hay coincidencia &
    $\log n + k$ &
    $O(\log n)$\newline$+ O(k)$ &
    $O(1)$, sin estructuras auxiliares \\
\addlinespace

registrarPosicionEnMapa &
    Inserción al final de la lista de recién ingresados &
    $1$ operación\newline(amortizado) &
    $O(1)$\newline amortizado &
    $O(1)$ \\
\addlinespace

fusionarRecienIngresados &
    Mezcla ordenada de zona ordenada ($n$) y lista de recién
    ingresados ($k$) &
    $(n+k)\log(n+k)$ &
    $O(n \log n)$ &
    $O(n+k)$, arreglo temporal para la mezcla \\
\addlinespace

buscarPosicionesPorCampo &
    Binaria en zona ordenada, lineal en lista de recién ingresados &
    $\log n + k$ &
    $O(\log n)$\newline$+ O(k)$ &
    $O(1)$ \\
\addlinespace

construirMapaPorCampo\newline(apellido) &
    Copia de $n$ llaves más ordenamiento &
    $n + n\log n$ &
    $O(n \log n)$ &
    $O(n)$, arreglo nuevo del tamaño del mapa \\

\bottomrule
\end{longtable}
}

La fila de \texttt{fusionarRecienIngresados} ilustra por qué reportar
solo $O(n \log n)$ no basta: el costo real es $(n+k)\log(n+k)$, y aunque
$k$ es pequeño frente a $n$ (acotado al 10\% de la zona ordenada, ADR
001), sigue sumando trabajo real en cada reorganización. No desaparece
por ser un término de menor orden.

\subsection{Módulo archivo (persistencia)}

{\small
\begin{longtable}{p{3.3cm} p{4.3cm} p{2.6cm} p{1.8cm} p{2.9cm}}
\toprule
Función & Algoritmo & Costo real & Big O & Costo espacial \\
\midrule
\endhead

leerRegistroEnOffset &
    \texttt{fseek} más \texttt{fread} de un registro de tamaño fijo &
    $1$ operación\newline de E/S &
    $O(1)$ &
    $O(1)$, un solo struct en RAM \\
\addlinespace

escribirRegistroEnOffset &
    \texttt{fseek} más \texttt{fwrite} de un registro de tamaño fijo &
    $1$ operación\newline de E/S &
    $O(1)$ &
    $O(1)$ \\
\addlinespace

purgarInactivos &
    Recorrido secuencial de $n$ registros, reescritura de los activos
    ($a \leq n$) y reconstrucción de 3 mapas &
    $n + a + n\log n$ &
    $O(n \log n)$ &
    $O(n)$, archivo temporal más mapas reconstruidos \\

\bottomrule
\end{longtable}
}

\subsection{Módulo estudiante (lógica de negocio)}

{\small
\begin{longtable}{p{3.3cm} p{4.3cm} p{2.6cm} p{1.8cm} p{2.9cm}}
\toprule
Función & Algoritmo & Costo real & Big O & Costo espacial \\
\midrule
\endhead

registrarEstudiante &
    Inserción en el mapa más escritura de 3 registros de tamaño fijo &
    $1 + 3$ &
    $O(1)$\newline amortizado &
    $O(1)$ \\
\addlinespace

buscarEstudiantePorCarne &
    Búsqueda en el mapa más 1 lectura en offset &
    $\log n + 1$ &
    $O(\log n)$ &
    $O(1)$ \\
\addlinespace

buscarEstudiantesPorApellido &
    Búsqueda en el mapa de apellido más $k$ lecturas de coincidencias &
    $\log n + k$ &
    $O(\log n)$\newline$+ O(k)$ &
    $O(k)$, lista de resultados \\
\addlinespace

modificarEstudiante &
    Búsqueda en el mapa más escritura en offset existente &
    $\log n + 1$ &
    $O(\log n)$ &
    $O(1)$ \\
\addlinespace

eliminarEstudianteLogico &
    Búsqueda en el mapa más marcar \texttt{estado = 0} en mapa y
    archivo &
    $\log n + 2$ &
    $O(\log n)$ &
    $O(1)$ \\
\addlinespace

consultarExpedienteCompleto &
    Búsqueda en el mapa más 3 lecturas (identidad, académico,
    contacto) &
    $\log n + 3$ &
    $O(\log n)$ &
    $O(1)$ \\
\addlinespace

mostrarTodosLosEstudiantes &
    Recorrido secuencial completo &
    $n$ &
    $O(n)$ &
    $O(1)$ en streaming,\newline $O(n)$ si se acumula \\
\addlinespace

mostrarEstudiantesPorCarrera &
    Recorrido secuencial con filtro &
    $n$ comparaciones &
    $O(n)$ &
    $O(r)$, $r$ resultados de esa carrera \\
\addlinespace

actualizarPromedio\newline/ actualizarCreditosAprobados &
    Búsqueda en el mapa más escritura en \texttt{academico.dat} &
    $\log n + 1$ &
    $O(\log n)$ &
    $O(1)$ \\

\bottomrule
\end{longtable}
}

No se justifica un mapa propio para \texttt{mostrarEstudiantesPorCarrera}:
con aproximadamente 50 carreras, un mapa adicional pagaría su propio
costo de mantenimiento (inserción, lista de recién ingresados,
reorganización) por una operación que ya es $O(n)$ y donde $n$ son los
mismos 25\,000 a un millón de registros que cualquier reporte necesita
tocar de todas formas.

\subsection{Módulo reportes (agregados y mantenimiento)}

{\small
\begin{longtable}{p{3.3cm} p{4.3cm} p{2.6cm} p{1.8cm} p{2.9cm}}
\toprule
Función & Algoritmo & Costo real & Big O & Costo espacial \\
\midrule
\endhead

calcularPromedioInstitucional &
    Recorrido secuencial con acumulador &
    $n$ sumas\newline más $1$ división &
    $O(n)$ &
    $O(1)$, un acumulador \\
\addlinespace

encontrarMayorPromedio &
    Recorrido secuencial con máximo &
    $n$ comparaciones &
    $O(n)$ &
    $O(1)$ \\
\addlinespace

encontrarMayorCreditos &
    Recorrido secuencial con máximo &
    $n$ comparaciones &
    $O(n)$ &
    $O(1)$ \\
\addlinespace

generarReporteEstadistico &
    Recorrido secuencial con múltiples acumuladores (por carrera, por
    nivel, etc.) &
    $n \times c$,\newline $c$ = cantidad de\newline categorías &
    $O(n)$ &
    $O(c)$, un acumulador por categoría \\

\bottomrule
\end{longtable}
}

El costo real de \texttt{generarReporteEstadistico} es $n \times c$ y no
simplemente $n$: por cada uno de los $n$ registros se actualizan $c$
acumuladores (total por carrera, por nivel académico, etc.). Como $c$ es
una constante pequeña y fija que no crece con $n$, la clase asintótica se
mantiene en $O(n)$, pero el costo real ejecutado es mayor que un
recorrido con un solo acumulador.
